\subsection{CLIP Text Encoder for Semantic Extraction}
\label{sec:method_clip_text_encoder}

To bridge the gap between human-understandable fault descriptions and the machine-learnable vibration manifold, we leverage the pre-trained CLIP (Contrastive Language-Image pretraining) text encoder as our primary semantic extractor. Unlike traditional zero-shot learning (ZSL) frameworks that rely on binary attribute vectors—which are often sparse, noisy, and labor-intensive to annotate—our approach utilizes the rich, continuous latent space of natural language to capture fine-grained diagnostic knowledge. This shift from binary to natural language semantics allows for a more expressive representation of industrial fault modes.

The process begins with the deliberate design of a structured fault description template. For each fault category $c \in \mathcal{S} \cup \mathcal{U}$, we construct a descriptive prompt $\mathcal{P}_c$ that encapsulates the physical characteristics, diagnostic symptoms, and spectral signatures of the specific fault mode. A typical template is formulated as follows: \textit{``A [component] fault characterized by [symptom 1], [symptom 2], and exhibiting [spectral feature] during steady-state operation.''} For instance, an inner race bearing fault might be described as \textit{``A bearing fault characterized by periodic impulsive shocks, increased vibration amplitude at the ball pass frequency, and metallic clicking sounds, often resulting from surface fatigue on the inner raceway.''} This templated approach ensures a consistent syntactic structure across different fault types while allowing for the inclusion of highly technical domain knowledge that binary attributes fail to capture.

The designed text prompt $\mathcal{P}_c$ is then fed into the frozen CLIP text encoder, denoted as $\mathcal{E}_{text}(\cdot)$. We adopt the Transformer-based architecture of the CLIP model (specifically the ViT-B/32 variant), which consists of a robust multi-head self-attention mechanism capable of capturing complex semantic dependencies between technical diagnostic terms. During our training phase, the weights of $\mathcal{E}_{text}$ remain strictly frozen to preserve the broad, cross-modal semantic knowledge learned during its large-scale pretraining on hundreds of millions of text-image pairs. The primary output of the encoder for a given class $c$ is a continuous semantic vector $\mathbf{s}_c \in \mathbb{R}^d$:
\begin{equation}
    \mathbf{s}_c = \mathcal{E}_{text}(\text{tokenize}(\mathcal{P}_c)),
\end{equation}
where $d$ (typically 512 for the ViT-B/32 variant) represents the dimensionality of the CLIP latent space.

This continuous semantic representation offers two significant advantages over standard binary attribute descriptors. First, it provides a much higher representational density and richness. The relative positions and cosine distances between vectors in $\mathbb{R}^d$ inherently reflect the semantic similarity between different fault modes; for example, different types of rolling element bearing faults will naturally cluster closer to each other in the latent space than they would to a rotor misalignment or a gear pitting fault. Second, the use of natural language offers unprecedented flexibility and scalability. Diagnostic experts can refine, expand, or augment the semantic descriptions using natural language without the need to retrain the entire feature extraction pipeline. This allows for a more intuitive and modular knowledge transfer mechanism that is essential for complex, evolving industrial systems, effectively transforming the zero-shot problem into a cross-modal alignment task.
